\section{Discussion and Limitations}\label{sec:discussion}

\paragraph{Novelty, not malice.} The model answers whether a container did something it never did. Two consequences
follow from our measurements. First, an attack that changes only the meaning of data inside a process is invisible to
it: the Juice-Shop SQL injections that LID-DS records are login requests to the same endpoint as normal logins, with a
longer body, no new process, path, error or size, and the traces keep only the first ${\sim}80$ bytes of each buffer.
No detector that reads system calls alone can see them; an application-layer monitor (a web application firewall,
runtime application self-protection, or query monitoring in the database driver) is the right complement. We did not
add a signal for this exploit after seeing it missed, which would be tuning on the test data. Second, a harmless new
process looks new too. The kernel lineage tag separates the workload's own process tree from operator sessions; whether
the native stream alone stays quiet on operator activity on a live cluster is the question of our latest run
\pending{H1}.

\paragraph{Adaptive attackers.} Every detection on LID-DS~2021 rests on sequences the workload never produced: an
attacker who removes all of them is caught 0/544 times in our semi-synthetic test, by the model and by the simpler
count rule alike. This matches the classic mimicry results against sequence-based detectors~\cite{wagner2002,tan2002}
and the more recent ones against provenance-based detectors~\cite{goyal2023}. Return values raise the cost of mimicry,
because the attacker must also reproduce sizes, paths and error codes, but they do not remove it. A rate signal over seen
sequences did not add detections on the datasets; its value against an attacker who pads an exploit with frequent seen
calls is what the dilution trials of the latest run measure \pending{H3}. Two further evasion paths are structural: a
load above about 400{,}000 new grams per second per cgroup fills the gram map, and the worker cannot tell an incomplete
window from a complete one; and a living-off-the-land attack with tools shipped in the image is exactly what the R1
rule would downgrade \pending{H4}.

\paragraph{Thresholds and changing normal behaviour.} The Gumbel assumption fails on both the datasets and the
cluster; held-out exceedances are 10 to 17 times the design level, so $\alpha'$ is a knob whose false-alarm rate must be
measured. Normal behaviour also changes with load, versions and incidents: after a cluster change, unchanged models put
2.3 times more blocks above threshold. Periodic refits are the natural answer, with a risk in the other direction: if an
attack enters the learning data, the model learns its sequences. Our fits exclude the spans around every trial because
we know the campaign schedule; an operational system would have to exclude alarms confirmed as attacks before refitting.

\paragraph{Scope of the evidence.} The live evidence comes from one cluster running one application, with self-written
attack fixtures and labels given by the trial design rather than by independent review. All campaign attacks enter
through \texttt{kubectl exec}; we did not run exploits of vulnerable software inside services on the cluster, so the
native stream is tested against attacks only offline, on LID-DS, where every event is native. The bridge has no durable
queue, and lost 113 of 190 alert episodes during a 3.5-minute host stall. RCA quality is measured by coverage and the
rank of the attack binary on one attack type, not by investigators. The overhead was measured with microbenchmarks, not
per request of a real application.
